\section{Proofs for Section~\ref{sec:local-rates}}\label{sec:local-proofs}

\subsection{Proof of Theorem~\ref{thm:composite-expansion}}\label{sec:local-proofs-composite}

Fix the compact set $\mathcal K$ of shapes and let
$\mathcal S_{\mathcal K}=\{(d,\xi):d\in\mathcal K,\ \xi\in D(d)\}$, a compact
set. Let $R_{\mathcal K}=\max_{d\in\mathcal K}\norm d_\infty$, so that
$\norm\xi_\infty\le R_{\mathcal K}$ on $\mathcal S_{\mathcal K}$ and
$B_t\subseteq x^*+tR_{\mathcal K}[-1,1]^n$. A quantity depending on
$(t,d,\xi)$ is $O(t^j)$ if it is bounded by $Ct^j$ on
$\mathcal S_{\mathcal K}$ for $0<t\le t_1$, with $C$ and $t_1$ depending
only on $\mathcal K$ and the factorable representation, and it is $o(t^2)$
if its supremum over $\mathcal S_{\mathcal K}$ divided by $t^2$ tends to
zero. For a factor index $a$ we abbreviate $a=v_a(x_t)$, $a^*=v_a^*$,
$a^L=v_a^L(B_t)$, $a^U=v_a^U(B_t)$, and write $\mathrm{cv}_a,\mathrm{cc}_a$
for the relaxations at $x_t$ on $B_t$; the letter $b$ is used in the same
way for a second factor. We first collect five lemmas.

\begin{lemma}[Median rule]\label{lem:mid}
Let $\mathrm{cv}\le v\le\mathrm{cc}$, $v\in Z=[L,U]$, and $z\in Z$. Then
$m=\midop\{\mathrm{cv},\mathrm{cc},z\}$ lies in $[\mathrm{cv},\mathrm{cc}]\cap Z$.
If $z=L$, then $m=\max\{\mathrm{cv},L\}$; if $z=U$, then
$m=\min\{\mathrm{cc},U\}$.
\end{lemma}

\begin{proof}
Since $\mathrm{cv}\le\mathrm{cc}$, $m$ is the projection of $z$ onto
$[\mathrm{cv},\mathrm{cc}]$. If $z<\mathrm{cv}$, then $m=\mathrm{cv}$ and
$L\le z<\mathrm{cv}\le v\le U$. If $z>\mathrm{cc}$, then $m=\mathrm{cc}$ and
$L\le v\le\mathrm{cc}<z\le U$. Otherwise $m=z\in Z$. If $z=L$, then
$L\le v\le\mathrm{cc}$, and the projection of $L$ onto $[\mathrm{cv},\mathrm{cc}]$
is $\max\{\mathrm{cv},L\}$. The case $z=U$ is symmetric.
\end{proof}

\begin{lemma}[Lipschitz envelopes]\label{lem:lipschitz-envelope}
If $\psi$ is $\Lambda$-Lipschitz on $[\alpha,\beta]$, its convex and concave
envelopes on $[\alpha,\beta]$ are $\Lambda$-Lipschitz.
\end{lemma}

\begin{proof}
Let $c$ be the convex envelope. The affine functions
$A_\beta(y)=\psi(\beta)-\Lambda(\beta-y)$ and
$A_\alpha(y)=\psi(\alpha)-\Lambda(y-\alpha)$ are minorants of $\psi$, so
$c\ge A_\beta$ and $c\ge A_\alpha$. Since $c\le\psi$ and
$A_\beta(\beta)=\psi(\beta)$, we get $c(\beta)=\psi(\beta)$, and likewise
$c(\alpha)=\psi(\alpha)$. For $\alpha\le y<y'<\beta$, convexity gives
\[
  \frac{c(y')-c(y)}{y'-y}\le\frac{c(\beta)-c(y)}{\beta-y}
  \le\frac{\psi(\beta)-A_\beta(y)}{\beta-y}=\Lambda,
\]
and for $\alpha<y<y'\le\beta$,
\[
  \frac{c(y')-c(y)}{y'-y}\ge\frac{c(y)-c(\alpha)}{y-\alpha}
  \ge\frac{A_\alpha(y)-\psi(\alpha)}{y-\alpha}=-\Lambda.
\]
The remaining cases follow by continuity of $c$ at the endpoints, which
holds because $A_\alpha\le c\le\psi$ near $\alpha$ and $A_\beta\le c\le\psi$
near $\beta$. For the concave envelope apply this to $-\psi$.
\end{proof}

\begin{lemma}[Envelopes of a $C^2$ function on a short interval]\label{lem:envelope-gap}
Let $h$ be twice continuously differentiable on $[a^*-r,a^*+r]$, let
$c=h''(a^*)$, and let $\omega(s)=\sup\{\abs{h''(y)-c}:\abs{y-a^*}\le s\}$,
so that $\omega(s)\to0$ as $s\downarrow0$, and $\omega(s)\le M_2s$ if $h''$ is
$M_2$-Lipschitz. For every interval $Z=[L,U]\subseteq[a^*-s,a^*+s]$ with
$s\le r$, and every $y\in Z$,
\begin{align*}
  \bigl|h(y)-h^{\mathrm{cv}}_Z(y)-\tfrac12c^-(y-L)(U-y)\bigr|&\le\omega(s)s^2,\\
  \bigl|h^{\mathrm{cc}}_Z(y)-h(y)-\tfrac12c^+(y-L)(U-y)\bigr|&\le\omega(s)s^2.
\end{align*}
Moreover $h^{\mathrm{cv}}_Z-h'(a^*)\,\mathrm{id}$ and
$h^{\mathrm{cc}}_Z-h'(a^*)\,\mathrm{id}$ are Lipschitz on $Z$ with constant
$s\sup_{\abs{y-a^*}\le s}\abs{h''(y)}$.
\end{lemma}

\begin{proof}
Let $q(y)=h(a^*)+h'(a^*)(y-a^*)+\tfrac c2(y-a^*)^2$ and $e=h-q$. Then
$e(a^*)=e'(a^*)=0$ and $\abs{e''}\le\omega(s)$ on $[a^*-s,a^*+s]$, so
$\abs e\le\eta=\omega(s)s^2/2$ on $Z$. Envelopes are monotone and commute
with adding constants, so $q^{\mathrm{cv}}_Z-\eta\le h^{\mathrm{cv}}_Z\le
q^{\mathrm{cv}}_Z+\eta$, and therefore
$\abs{(h-h^{\mathrm{cv}}_Z)-(q-q^{\mathrm{cv}}_Z)}\le2\eta$. If $c\ge0$, $q$ is
convex and $q-q^{\mathrm{cv}}_Z=0$. If $c<0$, $q$ is concave; its convex
envelope on $Z$ is the chord through its endpoint values, because the chord
is convex and below $q$, and every convex minorant of $q$ lies below the
chord. The difference $q-\text{chord}$ is a quadratic with leading
coefficient $c/2$ that vanishes at $L$ and $U$, so it equals
$\tfrac12\abs c(y-L)(U-y)$. The concave envelope is analogous. For the
Lipschitz claim, envelopes commute with adding affine functions, so
$h^{\mathrm{cv}}_Z-h'(a^*)\,\mathrm{id}$ is the convex envelope of
$\psi=h-h'(a^*)\,\mathrm{id}$, and $\abs{\psi'}\le s\sup\abs{h''}$ on $Z$;
Lemma~\ref{lem:lipschitz-envelope} applies.
\end{proof}

\begin{lemma}[Monotone functions]\label{lem:monotone-argmin}
If $h'>0$ on $Z=[L,U]$, then $L$ is the unique minimizer of $h^{\mathrm{cv}}_Z$
and $U$ is the unique maximizer of $h^{\mathrm{cc}}_Z$ on $Z$. If $h'<0$ on
$Z$, the roles of $L$ and $U$ are exchanged.
\end{lemma}

\begin{proof}
Let $\mu=\min_Zh'>0$. The affine function $h(L)+\mu(y-L)$ is a convex
minorant of $h$ on $Z$, so for $y>L$,
$h^{\mathrm{cv}}_Z(y)\ge h(L)+\mu(y-L)>h(L)\ge h^{\mathrm{cv}}_Z(L)$. Similarly
$h(U)+\mu(y-U)$ is a concave majorant, which gives the claim about
$h^{\mathrm{cc}}_Z$. The case $h'<0$ is symmetric.
\end{proof}

\begin{lemma}[Sign selections]\label{lem:sign-selection}
Let $\abs{s_t-s^*}\le Ct$ and $\mathrm{cv}\le v\le\mathrm{cc}$ with
$\mathrm{cc}-\mathrm{cv}\le C't^2$. Let $\tilde v=\mathrm{cv}$ if $s^*>0$,
$\tilde v=\mathrm{cc}$ if $s^*<0$, and $\tilde v=v$ if $s^*=0$. Then
$\abs{\min\{s_t\mathrm{cv},s_t\mathrm{cc}\}-s_t\tilde v}\le CC't^3$, and the
difference is zero if $s^*\ne0$ and $t<\abs{s^*}/C$. The same holds for the
maximum, with $\mathrm{cv}$ and $\mathrm{cc}$ exchanged in the definition of
$\tilde v$.
\end{lemma}

\begin{proof}
The minimum equals $s_t\mathrm{cv}$ if $s_t\ge0$ and $s_t\mathrm{cc}$ if
$s_t\le0$. If $s_t=0$, or if $s^*\ne0$ and $s_t$ has the sign of $s^*$, the
difference is zero; this holds in particular when $s^*\ne0$ and
$t<\abs{s^*}/C$. In the remaining cases, $s^*=0$ or $s_t$ and $s^*$ have
opposite signs, so $\abs{s_t}\le\abs{s_t-s^*}\le Ct$; both the minimum and
$s_t\tilde v$ lie between $s_t\mathrm{cv}$ and $s_t\mathrm{cc}$, an interval of
length $\abs{s_t}(\mathrm{cc}-\mathrm{cv})\le CC't^3$.
\end{proof}

\paragraph{Proof of the theorem.}
By induction over $k$, each $v_k$ is twice continuously differentiable on a
neighborhood $N$ of $x^*$: sums, scalings, and products preserve this, and
$h_k(v_a)$ is twice continuously differentiable near $x^*$ because $v_a$ is
continuous and $h_k$ is twice continuously differentiable near $v_a(x^*)$.
For $t$ small, $x_t\in N$, and Taylor's theorem gives, uniformly on
$\mathcal S_{\mathcal K}$,
\begin{equation}\label{eq:taylor-factor}
  v_k(x_t)=v_k^*+t\pi_k(\xi)+O(t^2).
\end{equation}
We prove~\eqref{eq:expansion-interval}--\eqref{eq:expansion-dominance} by
induction over $k$; we refer to them as $(\mathrm I_k)$, $(\mathrm R_k)$, and
$(\mathrm P_k)$. By $(\mathrm I_a)$, the interval of a univariate argument
lies in $[a^*-Ct,a^*+Ct]$, so for small $t$ it lies in the interval where
$h_k$ is twice continuously differentiable, and all objects are defined.
By $(\mathrm R_a)$ and boundedness of $E_a$ on the compact set
$\mathcal S_{\mathcal K}$, $\mathrm{cc}_a-\mathrm{cv}_a=O(t^2)$.

\emph{Variables and constants.} The intervals are exact, and
$\mathrm{cv}=\mathrm{cc}=v$, so all three properties hold with $E=0$.

\emph{Sums and scalings.} The expansions add. For a sum,
$v^L-\mathrm{cv}=(a^L-\mathrm{cv}_a)+(b^L-\mathrm{cv}_b)\le o(t^2)$. For $c<0$,
$\mathrm{cv}=c\,\mathrm{cc}_a=ca-t^2\abs cE_a^{\mathrm{cc}}+o(t^2)$ and
$v^L-\mathrm{cv}=ca^U-c\,\mathrm{cc}_a=\abs c(\mathrm{cc}_a-a^U)\le o(t^2)$.

\emph{Products $v_k=v_av_b$: intervals.} Each endpoint product is
$a^ib^j=a^*b^*+t(b^*\pi_a^i+a^*\pi_b^j)+O(t^2)$, where $\pi_a^i$ ranges over
$\underline\pi_a,\overline\pi_a$ and $\pi_b^j$ over
$\underline\pi_b,\overline\pi_b$. The minimum and maximum are
$1$-Lipschitz in the sup norm, and
$\min_{i,j}(b^*\pi_a^i+a^*\pi_b^j)=\min_ib^*\pi_a^i+\min_ja^*\pi_b^j$, which
is the lower endpoint of the interval in the recursion; the maximum is
analogous. This proves $(\mathrm I_k)$.

\emph{Products: relaxations.} By $(\mathrm I_a)$ and $(\mathrm I_b)$, the
coefficients $b^L,b^U$ are within $O(t)$ of $b^*$ and $a^L,a^U$ are within
$O(t)$ of $a^*$. Lemma~\ref{lem:sign-selection} gives
\[
  \alpha_1=b^L\tilde a+O(t^3),\quad\beta_1=b^U\tilde a+O(t^3),\quad
  \alpha_2=a^L\tilde b+O(t^3),\quad\beta_2=a^U\tilde b+O(t^3),
\]
with the same $\tilde a\in\{\mathrm{cv}_a,\mathrm{cc}_a,a\}$, selected by the
sign of $b^*$, in $\alpha_1$ and $\beta_1$, and the same $\tilde b$, selected
by the sign of $a^*$, in $\alpha_2$ and $\beta_2$. This is the key point: for
small $t$, $b^L$ and $b^U$ have the sign of $b^*$ when $b^*\ne0$, and the
selection costs only $O(t^3)$ when $b^*=0$. The identities
\[
  b^L\alpha+a^L\beta-a^Lb^L=\alpha\beta-(\alpha-a^L)(\beta-b^L),\qquad
  b^U\alpha+a^U\beta-a^Ub^U=\alpha\beta-(a^U-\alpha)(b^U-\beta)
\]
and the fact that the maximum is $1$-Lipschitz give
\begin{equation}\label{eq:product-mccormick}
  \mathrm{cv}_k=\tilde a\tilde b
  -\min\bigl\{(\tilde a-a^L)(\tilde b-b^L),\ (a^U-\tilde a)(b^U-\tilde b)\bigr\}
  +O(t^3).
\end{equation}
Since $\tilde a-a=O(t^2)$, $\tilde b-b=O(t^2)$, $b-b^*=O(t)$, and $a-a^*=O(t)$,
\[
  \tilde a\tilde b=ab+b^*(\tilde a-a)+a^*(\tilde b-b)+O(t^3).
\]
By $(\mathrm R_a)$, $b^*(\tilde a-a)=-t^2\bigl((b^*)^+E_a^{\mathrm{cv}}+(b^*)^-E_a^{\mathrm{cc}}\bigr)+o(t^2)$:
the two cases are $\tilde a=\mathrm{cv}_a$ with $b^*>0$ and
$\tilde a=\mathrm{cc}_a$ with $b^*<0$. The term $a^*(\tilde b-b)$ is
analogous. Next, by~\eqref{eq:taylor-factor} and $(\mathrm I_a)$,
$\tilde a-a^L=(a-a^L)+(\tilde a-a)=t(\pi_a-\underline\pi_a)+O(t^2)$, and
likewise for the other three differences, so each product in the minimum
of~\eqref{eq:product-mccormick} is $t^2$ times the corresponding product in
$E_k^{\mathrm{cv}}$, plus $O(t^3)$. This proves the convex half of
$(\mathrm R_k)$. The concave half is the same argument with the maximum
selections, which select $\mathrm{cc}_a$ for $b^*>0$ and $\mathrm{cv}_a$ for
$b^*<0$, and the identities
$b^L\alpha+a^U\beta-a^Ub^L=\alpha\beta+(a^U-\alpha)(\beta-b^L)$ and
$b^U\alpha+a^L\beta-a^Lb^U=\alpha\beta+(\alpha-a^L)(b^U-\beta)$.

\emph{Products: dominance.} The function
$M(\alpha,\beta)=\max\{b^L\alpha+a^L\beta-a^Lb^L,\ b^U\alpha+a^U\beta-a^Ub^U\}$
is, on $\Pi=[a^L,a^U]\times[b^L,b^U]$, the convex envelope of $\alpha\beta$,
so it is at least $\min_\Pi\alpha\beta=v_k^L$ there. On $\R^2$ it is Lipschitz in
the sup norm with constant $\max\{\abs{a^L}+\abs{b^L},\abs{a^U}+\abs{b^U}\}=O(1)$.
Since $\tilde a\in[\mathrm{cv}_a,\mathrm{cc}_a]$, $(\mathrm P_a)$ gives
$a^L-o(t^2)\le\tilde a\le a^U+o(t^2)$, and similarly for $\tilde b$, so
$(\tilde a,\tilde b)$ lies within sup-distance $o(t^2)$ of $\Pi$. Together
with~\eqref{eq:product-mccormick}, which states $\mathrm{cv}_k=M(\tilde a,\tilde b)+O(t^3)$,
this gives $v_k^L-\mathrm{cv}_k\le o(t^2)$. The concave side is symmetric.

\emph{Univariate factors $v_k=h(v_a)$.} Let $Z=[a^L,a^U]$, which lies in
$[a^*-s_t,a^*+s_t]$ with $s_t=Ct$, and let $h_1=h'(a^*)$, $h_2=h''(a^*)$.

For $(\mathrm I_k)$: on $Z$, $h(y)=h(a^*)+h_1(y-a^*)+O(t^2)$, and the minimum
over $Z$ of $h_1(y-a^*)$ is $t\min\{h_1\underline\pi_a,h_1\overline\pi_a\}+O(t^2)$.

For $(\mathrm P_k)$, which here holds exactly: let
$m=\midop\{\mathrm{cv}_a,\mathrm{cc}_a,z^{\min}\}$. By Lemma~\ref{lem:mid},
$m\in Z$, so $\mathrm{cv}_k=h^{\mathrm{cv}}_Z(m)\ge\min_Zh^{\mathrm{cv}}_Z=\min_Zh=v_k^L$;
the convex envelope has the same minimum as $h$, because the constant
$\min_Zh$ is a convex minorant. The concave side is analogous.

For $(\mathrm R_k)$ write
\[
  \mathrm{cv}_k=h(a)-\bigl[h(a)-h^{\mathrm{cv}}_Z(a)\bigr]
  +\bigl[h^{\mathrm{cv}}_Z(m)-h^{\mathrm{cv}}_Z(a)\bigr].
\]
By Lemma~\ref{lem:envelope-gap}, the first bracket equals
$\tfrac12h_2^-(a-a^L)(a^U-a)+O(\omega(s_t)s_t^2)=\tfrac12t^2h_2^-\sigma_a+o(t^2)$.
This holds in every case, including $h_2=0$ and sign changes of $h''$ inside
$Z$, because the comparison is with the quadratic model of $h$ and not with
the convexity of $h$ on $Z$. By Lemma~\ref{lem:mid}, $m,a\in Z$ and
$\abs{m-a}\le\mathrm{cc}_a-\mathrm{cv}_a=O(t^2)$, so the Lipschitz part of
Lemma~\ref{lem:envelope-gap} gives
$h^{\mathrm{cv}}_Z(m)-h^{\mathrm{cv}}_Z(a)=h_1(m-a)+O(t)\,O(t^2)$. It remains to
evaluate $h_1(m-a)$.
\begin{itemize}
\item If $h_1=0$, the term is zero. The minimizer $z^{\min}$ may then lie
  anywhere in $Z$ and may jump as $t$ varies; this does not matter.
\item If $h_1>0$, then $h'>0$ on $Z$ for small $t$, Lemma~\ref{lem:monotone-argmin}
  gives $z^{\min}=a^L$, and Lemma~\ref{lem:mid} gives
  $m=\max\{\mathrm{cv}_a,a^L\}$. Hence
  $m-a=(\mathrm{cv}_a-a)+(a^L-\mathrm{cv}_a)^+$, and $(\mathrm P_a)$ gives
  $(a^L-\mathrm{cv}_a)^+=o(t^2)$, so $m-a=-t^2E_a^{\mathrm{cv}}+o(t^2)$. Without
  $(\mathrm P_a)$, the clipping of the argument at $a^L$ could change the
  second-order term; this is the reason for proving $(\mathrm P_k)$ along
  with $(\mathrm R_k)$.
\item If $h_1<0$, then $z^{\min}=a^U$, $m=\min\{\mathrm{cc}_a,a^U\}$, and
  $m-a=t^2E_a^{\mathrm{cc}}+o(t^2)$, so
  $h_1(m-a)=-t^2\abs{h_1}E_a^{\mathrm{cc}}+o(t^2)$.
\end{itemize}
Adding the three parts gives $\mathrm{cv}_k=v_k(x_t)-t^2E_k^{\mathrm{cv}}+o(t^2)$.
The concave side uses $z^{\max}$, which is $a^U$ if $h_1>0$ and $a^L$ if
$h_1<0$, and gives $E_k^{\mathrm{cc}}$. This completes the induction.

\emph{Remainder order.} If each $h_k''$ is Lipschitz near the relevant
point, then $\omega(s)=O(s)$ in Lemma~\ref{lem:envelope-gap}, every remainder
in the induction, including those in $(\mathrm P_k)$, is $O(t^3)$, and $f$ has
a Lipschitz second derivative near $x^*$, so its Taylor remainder is also
$O(t^3)$.

\emph{The objective.} Since $f=v_m$ is twice continuously differentiable near
$x^*$, $f(x_t)=f(x^*)+t\nabla f(x^*)^T\xi+\tfrac12t^2\xi^T\nabla^2f(x^*)\xi+o(t^2)$
uniformly; combine with $(\mathrm R_m)$. Continuity of $Q$ follows from
continuity of $E_m^{\mathrm{cv}}$.

\emph{Clipping.} Clipping does not change the intervals, which do not use
the relaxations. Clipped relaxations still bracket the factor values,
because $v_k^L\le v_k\le v_k^U$; hence Lemma~\ref{lem:mid} still applies.
Every rule is Lipschitz in its relaxation inputs with constants $O(1)$: the
product rule has interval endpoints as coefficients, and the univariate rule
composes the $1$-Lipschitz median with $h^{\mathrm{cv}}_Z$, which is
Lipschitz with constant $\max_Z\abs{h'}=O(1)$ by
Lemma~\ref{lem:lipschitz-envelope}. By induction, the rule applied to
clipped inputs differs from its unclipped value by $o(t^2)$, and the
clipping of its output moves it by at most
$(v_k^L-\mathrm{cv}_k)^++o(t^2)=o(t^2)$ by $(\mathrm P_k)$. Hence the clipped
relaxations have the same expansions.\qed

\subsection{Proofs for minimizers on the boundary}\label{sec:local-proofs-boundary}

\paragraph{Proof of Lemma~\ref{lem:boundary-step}.}
Let $\mathcal K=\{d\in\mathcal S:d_F\in\mathcal K_F,\ 0\le d_i^+\le1\text{ for }i\in A\}$,
a compact set. Let $\omega$ be a modulus of uniform continuity of $Q$ on the
compact set $\{(d,\xi):d\in\mathcal K,\ \xi\in D(d)\}$ with respect to the sup
norm, and let $\varrho=\varrho_{\mathcal K}$. Choose $\bar\sigma\in(0,1]$ with
$\omega(\bar\sigma)\le\delta/4$, and $\bar t$ such that $\varrho(t)\le\delta/4$
and the comparison boxes belong to $\mathfrak B$ for $t\le\bar t$. Let
$\sigma=\max_{i\in A}d_i^+\le\bar\sigma$. For $\xi\in D(d)$, the pair
$(\iota(d_F),\iota(\xi_F))$ is within sup-distance $\sigma$ of $(d,\xi)$,
because $d_A^-=0$ and $0\le\xi_A\le d_A^+$. Hence
\begin{equation}\label{eq:boundary-continuity}
  \abs{Q(d,\xi)-Q^F(d_F,\xi_F)}\le\omega(\sigma)\le\delta/4.
\end{equation}
A point $x=x^*+t\xi\in K_U(B)$ satisfies, by Assumption~\ref{ass:boundary},
\begin{equation}\label{eq:boundary-retained}
  tg^T\xi+t^2Q(d,\xi)\le\epsilon+\varrho(t)t^2,\qquad
  g^T\xi=\sum_{i\in A}g_i\xi_i\ge0.
\end{equation}

(i) Dropping $g^T\xi\ge0$ and using~\eqref{eq:boundary-continuity},
$Q^F(d_F,\xi_F)\le\epsilon/t^2+\varrho(t)+\omega(\sigma)\le\delta$, so $\xi_F$
lies in the sublevel set defining $\Phi^F_\delta(d_F)$.

(ii) For $\xi=\iota(\xi_F)$, $g^T\xi=0$ and $Q(d,\xi)\le-\delta+\delta/4$, so
$\phi_B(x^*+t\xi)\le f^*+t^2(-3\delta/4+\delta/4)<f^*\le U$.

(iii) By~\eqref{eq:boundary-retained} and~\eqref{eq:boundary-continuity},
for $i\in A$,
\[
  g_i\xi_i\le g^T\xi\le\frac{\epsilon}{t}+t\bigl(-Q(d,\xi)+\varrho(t)\bigr)
  \le\frac{\epsilon}{t}+t\bigl(G^F(d_F)+\omega(\sigma)+\varrho(t)\bigr).
\]
The lower bound of coordinate $i$ stays at $x_i^*$, because $x^*\in K_U(B)$
by validity. Hence
\begin{equation}\label{eq:boundary-active-exact}
  w_i(T_U(B))\le\bigl(\epsilon+(G^F(d_F)+\omega(\sigma)+\varrho(t))t^2\bigr)/g_i,
\end{equation}
which gives (iii).

(iv) If $\epsilon+(G^F(d_F)-\delta/2)t^2<0$, the asserted bound is negative and
holds trivially. Otherwise let $\xi_F$ minimize $Q^F(d_F,\cdot)$ over
$D_F(d_F)$, fix $i\in A$ and $0\le s\le d_i^+$, and let $\xi$ have free part
$\xi_F$, component $s$ in coordinate $i$, and zero in the other active
coordinates. Then
$\xi\in D(d)$ and, by~\eqref{eq:boundary-continuity},
\[
  \phi_B(x^*+t\xi)\le f^*+tg_is+t^2\bigl(-G^F(d_F)+\delta/4+\delta/4\bigr),
\]
which is at most $U$ whenever $g_is\le\epsilon/t+t(G^F(d_F)-\delta/2)$. The
right side is nonnegative, so $s=0$ qualifies; taking the largest such
$s\le d_i^+$ gives (iv).\qed

\paragraph{Proof of Theorem~\ref{thm:boundary}.}
Let $M'=\max\{M,1\}$ and
$\mathcal K'=\{d\in\mathcal S:d_F=u_F,\ 0\le d_i^+\le M'\text{ for }i\in A\}$.
Let $\omega$ be a modulus of uniform continuity of $Q$ on
$\{(d,\xi):d\in\mathcal K',\ \xi\in D(d)\}$, and $\varrho=\varrho_{\mathcal K'}$.
Choose $\bar\sigma\in(0,1]$ with $\omega(\bar\sigma)\le\delta/4$. Let
$C=(G_M+\delta)/(\lambda g_{\min})$, and choose $\bar t$ such that, for
$t\le\bar t$, $\varrho(t)\le\delta/4$, the comparison boxes belong to
$\mathfrak B$, and $C\bar t\le\bar\sigma$. The proof of
Lemma~\ref{lem:boundary-step} applies on $\mathcal K'$ with these constants,
including the exact form~\eqref{eq:boundary-active-exact}. By
Lemma~\ref{lem:sequential}, every bound below for Jacobi rounds also holds
for complete sequential rounds, since the result of such a round is
contained in the Jacobi image of the same box.

\emph{First round.} Let $C_0=x^*+t_0D(d^0)$. Then
$B_1=T_U(B_0)\subseteq T_U(C_0)$ by Lemma~\ref{lem:order}, and
$B_1\subseteq B_0$, so the free part of $B_1$ lies in
$x_F^*+t_1D_F(u_F)$; this is (a) for $k=1$. For $x=x^*+t_0\xi\in K_U(C_0)$,
\eqref{eq:boundary-retained} holds with $d=d^0$, and $-Q(d^0,\xi)\le G_M$ by
the definition of $G_M$. As in the proof of
Lemma~\ref{lem:boundary-step}(iii),
$w_i(B_1)\le(\epsilon+(G_M+\delta/4)t_0^2)/g_i$, the first claim of (b).

\emph{Induction.} Suppose $N\ge2$, so that $\epsilon\le\delta t_1^2/2$. Let
$\sigma_k=\max_{i\in A}w_i(B_k)/t_k$. Since $t_1=t_0$,
\[
  \sigma_1\le\frac{\delta/2+G_M+\delta/4}{g_{\min}}\,t_0\le C\bar t\le\bar\sigma.
\]
Suppose, for some $1\le k\le N-1$, that the free part of $B_k$ lies in
$x_F^*+t_kD_F(u_F)$ and $\sigma_k\le\bar\sigma$. Then
$B_k\subseteq C_k=x^*+t_kD(d^k)$, where $d^k_F=u_F$ and the active components
of $d^k$ are at most $\sigma_k$, so $d^k\in\mathcal K'$. Since $k<N$,
$t_k^2\ge2\epsilon/\delta$. Lemma~\ref{lem:boundary-step}(i) and
Lemma~\ref{lem:order} place the free part of $B_{k+1}$ in
$x_F^*+t_kD_F(\Phi^F_\delta(u_F))\subseteq x_F^*+t_{k+1}D_F(u_F)$,
and~\eqref{eq:boundary-active-exact} gives
\[
  w_i(B_{k+1})\le\bigl(\epsilon+(G_0+\eta_k)t_k^2\bigr)/g_i,\qquad
  \eta_k=\omega(\sigma_k)+\varrho(t_k)\le\delta/2.
\]
Using $\epsilon\le\delta t_k^2/2$ and $G_0\le G_M$,
\[
  \sigma_{k+1}\le\frac{G_0+\delta}{\lambda g_{\min}}\,t_k\le C\bar t\le\bar\sigma.
\]
This completes the induction and proves (a) for $k\le N$ and (b) for
$k\le N-1$. For $k=N$, the induction hypothesis holds at $B_N$, and
\eqref{eq:boundary-active-exact} does not require $\epsilon\le\delta t^2/2$, so
the estimate of (b) also holds for $k=N$. Finally $\sigma_1\le Ct_0$ and
$\sigma_k\le Ct_{k-1}$ for $2\le k\le N$, so $\eta_k=\omega(\sigma_k)+\varrho(t_k)$
tends to zero as $t_k\to0$, by the continuity and remainder moduli.

\emph{Limits.} If $\epsilon=0$, then $N=\infty$; the free parts shrink to
$x_F^*$ and the active widths are at most $(G_0+\eta_k)t_k^2/g_i\to0$, so
$B_\infty=\{x^*\}$, and $\eta_k\to0$ gives the stated asymptotic form. If
$\epsilon>0$, then $N<\infty$ and $B_\infty\subseteq B_N$, whose free part lies
in $x_F^*+t_ND_F(u_F)$ with $t_N<\sqrt{2\epsilon/\delta}$. For the active
widths, if $N\ge2$, then
\[
  w_i(B_\infty)\le w_i(B_{N+1})\le\Bigl(\epsilon+\bigl(G_0+\tfrac\delta2\bigr)\tfrac{2\epsilon}{\delta}\Bigr)\Big/g_i
  =\Bigl(2+\frac{2G_0}{\delta}\Bigr)\frac{\epsilon}{g_i},
\]
and if $N=1$, then $t_0^2=t_1^2<2\epsilon/\delta$ and
\[
  w_i(B_\infty)\le w_i(B_1)\le\Bigl(\epsilon+\bigl(G_M+\tfrac\delta4\bigr)\tfrac{2\epsilon}{\delta}\Bigr)\Big/g_i
  \le\Bigl(\frac32+\frac{2G_M}{\delta}\Bigr)\frac{\epsilon}{g_i}.
\]
Both are at most $(2+2G_M/\delta)\epsilon/g_i$. The last sentence of the
theorem follows from the definition of $r^*(\Phi^F)$.\qed
